\begin{abstract}
Son yıllarda makine öğrenmesindeki hâkim eğilim, zor problemleri
daha büyük modeller, daha güçlü hızlandırıcılar ve daha geniş bellek
bütçeleriyle çözmek oldu. Ancak yıllardır süren küresel yarı iletken
arz sıkıntısı~\cite{burkacky2022chipshortage,khan2021chipshortage}
ve sürekli çevrim içi çalışan çıkarım sistemlerinin artan enerji ve
karbon maliyeti~\cite{strubell2019energy,patterson2021carbon}, bu
yaklaşımın kırılganlığını açık biçimde ortaya koymaktadır. Bu tablo,
ters yönde bir arayışı da gerekli kılar: yapay zekâ ve makine
öğrenmesi algoritmalarını, giyilebilir cihazlarda, sensörlerde ve uç
sistemlerde zaten yaygın biçimde kullanılan küçük
mikrodenetleyicilere sığacak biçimde yeniden tasarlamak. Bu çalışma
bu arayışa odaklanmaktadır. Kaynak kısıtlı çıkarım için geliştirilmiş
kompakt bir kapılı yinelemeli hücre olan
\fastgrnn{}~\cite{kusupati2018fastgrnn} algoritmasını uçtan uca ve
açık kaynaklı olarak yeniden üretiyor; modeli, mevcut silikon
ekosisteminin alt ucunu temsil eden iki bare-metal hedefte
çalıştırıyoruz: 8-bit \arduino{} (ATmega328P) ve 16-bit \msp{}
(donanım çarpıcısı yok; 16~KB Flash; 512~B SRAM). Sıkıştırma
boru hattı düşük ranklı ağırlık ayrıştırmasını, iteratif sert
eşikleme ile seyreklik kazandırmayı ve açık aktivasyon kalibrasyonu
içeren tensör-başına Q15 eğitim-sonrası nicelendirmeyi bir araya
getirir. Ortaya çıkan modelin ağırlıkları yalnızca
\textbf{566~bayt} yer kaplar; model \hapt{} kıyaslama kümesinde
\textbf{makro F1 $=$ 0.918} değerine ulaşır ve 3{,}399 test
penceresinin tamamında PyTorch referansıyla \textbf{\%100 tahmin
uyumu} gösterir. Her iki platform da \textbf{50~Hz'de gerçek
zamanlı akış halinde çıkarım} yapabilir (Arduino'da örnek başına
9.21~ms; MSP430'da 13~ms). Ayrıca 256 girişli sigmoid/tanh arama
tablosu, donanım çarpıcısı bulunmayan MSP430 üzerinde
\textbf{30.5$\times$ hızlanma} sağlar. Çalışma özgün \fastgrnn{} makalesini dört noktada genişletir:
(i)~platformlar arasında bit-eşdeğer deterministik çıkarım;
(ii)~uçtan uca yanıt süresini belirleyen yaklaşık 2~saniyelik
yinelemeli ısınma gecikmesinin nicel incelenmesi;
(iii)~çarpıcısız gömülü hedeflerde \textbf{30.5$\times$} uçtan uca
hızlanma sağlayan uygulanabilir bir arama tablosu yöntemi; ve
(iv)~donanım akım sensörü (INA226) ile gerçekleştirilen enerji
karakterizasyonu: \textbf{17.7~mW} aktif çıkarım gücü,
\textbf{$<$0.09~mW} uyku gücü ve arama tablosuyla çıkarım penceresi
başına \textbf{\%96.7 enerji azalması}. Sonuçlar, kompakt yinelemeli
mimarilerin kalibre edilmiş nicelendirme, arama tablosu tabanlı
aktivasyonlar ve ölçülmüş enerji profili ile birlikte; doğru,
bellek-verimli ve enerji-verimli insan aktivitesi tanımayı özel
hızlandırıcılara gerek duymadan ultra-kısıtlı mikrodenetleyicilerde
mümkün kılabileceğini göstermektedir.
\end{abstract}
